\documentclass{article}

\usepackage[a4paper, left=1.5cm, right=1.5cm, top=2cm, bottom=2cm]{geometry}

\usepackage{../../../../components/components}
\usepackage{hyperref}
\usepackage{fancyhdr}


% Configuration des en-têtes et pieds de page
\pagestyle{fancy}
\fancyhf{} % reset tout

\fancyhead[L]{DL2 Math-Info ASF4}
\fancyhead[C]{Analyse}
\fancyhead[R]{2025-2026}

\fancyfoot[L]{Ewen Rodrigues de Oliveira}
\fancyfoot[R]{\thepage}

\begin{document}

\docTitle{Chapitre 9 : Intégrales doubles}

\warning{
    Cours non finalisé, non structuré et contenant beaucoup d'oublis.\\
    Il est pour l'instant peu sûr de réviser à partir de ce cours, mais il peut être intéressant de le parcourir pour se faire une idée du sujet.
}

\section{Définition et propriétés des intégrales doubles}

\subsection{Définition}

Soit $f : [a,b] \times [c,d] \to \mathbb{R}$. On veut définir $\iint_{[a,b] \times [c,d]} f(x,y) \, dx \, dy$.
Représentation géométrique :
\begin{itemize}
    \item l'intégrale simple évalue l'aire entre la courbe et l'axe des abscisses,
    \item l'intégrale double évalue le volume entre la courbe et les deux axes.
\end{itemize}

\reminder{
    $f$ est continue en $(x_0, y_0) \in [a,b] \times [c,d]$ 
    \begin{itemize}
        \item $\Leftrightarrow \; \lim_{x \to x_0, \; y \to y_0} f(x,y) = f(x_0, y_0)$,
        \item $\Leftrightarrow \; \forall \varepsilon > 0, \; \exists \delta > 0, \; \forall (x,y) \in [a,b] \times [c,d], \; |(x,y) - (x_0, y_0)| < \delta \Rightarrow |f(x,y) - f(x_0, y_0)| < \varepsilon$.
        \item $\Leftrightarrow$ pour toute suite $(x_n, y_n)$ de $[a,b] \times [c,d]$ qui converge vers $(x_0, y_0)$, la suite $f(x_n, y_n)$ converge vers $f(x_0, y_0)$.
    \end{itemize}
    Pour s'entrainer et avoir plus de formalisme sur la topologie sur $\mathbb{R}^n$, le cours \href{https://ewenrdo.fr/ressources}{est disponible ici} dans "Mathématiques > Licence 2 > Analyse 3 > Chapitre 4 : Topologie".
}

\theorem{Théorème de Fubini}{}{false}{
    Soit $f : [a,b] \times [c,d] \to \mathbb{R}$ continue. Alors :
    \[
    \iint_{[a,b] \times [c,d]} f(x,y) \, dx \, dy = \int_a^b \left( \int_c^d f(x,y) \, dy \right) dx = \int_c^d \left( \int_a^b f(x,y) \, dx \right) dy
    \]
}
\remark{L'ordre d'intégration n'a pas d'importance, ce qui est cohérent avec l'idée que l'intégrale double évalue un volume : je peux commencer par calculer l'aire de la section verticale (en intégrant par rapport à $y$) ou l'aire de la section horizontale (en intégrant par rapport à $x$), le résultat est le même.}


\example{
    Considérons la fonction définie sur $\mathbb{R}^2$ par :
    \[
        f(x,y) = \begin{cases}
            \frac{xy}{x^4 + y^4} & \text{si (x,y) $\neq$ (0,0)} \\
            0 & \text{si (x,y)=(0,0)}
        \end{cases}
    \]
    Si $x$ est fixé, on a : $f(0,y)= 0$ et $y \neq 0, \; f(x,y) = \frac{xy}{x^4 + y^4}$. Quel que soit $y$, les deux fonctions sont continues.\\
    De même, si $y$ est fixé, on a : $f(x,0)= 0$ et $x \neq 0, \; f(x,y) = \frac{xy}{x^4 + y^4}$. Quel que soit $x$, les deux fonctions sont continues.\\
    Cependant, $f$ n'est pas continue en $(0,0)$.\\
    Prenons en effet $(x_n, y_n) = \left( \frac{1}{n}, \frac{1}{n} \right)$, alors $f(x_n, y_n) = \frac{1/n^2}{2/n^4} = \frac{n^2}{2}$ qui ne converge pas vers $0$.
}

\example{
    Considérons la fonction définie sur $\mathbb{R}^2$ par :
    \[
        f(x,y) = \begin{cases}
            \frac{xy}{\sqrt{x^2 + y^2}} & \text{si (x,y) $\neq$ (0,0)} \\
            0 & \text{si (x,y)=(0,0)}
        \end{cases}
    \]
    Posons $f(x_n, y_n) = \frac{x_n y_n}{\sqrt{x_n^2 + y_n^2}}$. On sait que $|x_n y_n| \leq \frac{x_n^2 + y_n^2}{2}$.\\
    Donc $0 \leq |f(x_n, y_n)| \leq \frac{\sqrt{x_n^2 + y_n^2}}{2}$ qui converge vers $0$.\\ Donc $f$ est continue en $(0,0)$ et donc continue sur $\mathbb{R}^2$.
}
\remark{Ce nouveau type d'intégral est particulièrement intéressant, notamment lorsqu'on s'intéresse au domaine $D$ d'intégration qui n'est pas un rectangle, mais une région plus générale du plan.}
\example{
    \begin{itemize}
        \item $D = \{ (x,y) \in \mathbb{R}^2 \; | \; x^2 + y^2 \leq 1\}$ est le disque de centre $0$ et de rayon $1$,
        \item $D = \{ (x,y) \in \mathbb{R}^2 \; | \frac{1}{2}x = x^2+y^2\}$ est un donut entre les cercles de rayon $\frac{1}{2}$ et $1$,
    \end{itemize}    
}

\theorem{Théorème}{}{false}{
    Soit $D \subset \mathbb{R}^2$ un sous-ensemble fermé borné quarrable.\\
    Soit $f : D \to \mathbb{R}$ une fonction continue.\\
    Alors, on peut définir l'intégrale double $\iint_D f \, dx \, dy$ telle que :
    \begin{enumerate}
        \item \textbf{Linéarité :} $\iint_D (\alpha f(x,y) + \beta g(x,y)) \, dx \, dy = \alpha \iint_D f(x,y) \, dx \, dy + \beta \iint_D g(x,y) \, dx \, dy$ (où $g : D \to \mathbb{R}$ est une autre fonction continue et $\alpha, \beta \in \mathbb{R}$),
        \item \textbf{Positivité :} si $f$ est positive sur $D$, alors $\iint_D f(x,y) \, dx \, dy \geq 0$,
        \item \textbf{Analogue de la relation de Chasles :} si $D = \bigsqcup_{i=1}^n D_i$ où les $D_i$ sont des sous-ensembles fermés bornés quarrables de $\mathbb{R}^2$, alors $\iint_D f(x,y) \, dx \, dy = \sum_{i=1}^n \iint_{D_i} f(x,y) \, dx \, dy$,
        \item Si $f \geq 0$, et $D \subset D'$, alors $\iint_D f(x,y) \, dx \, dy \leq \iint_{D'} f(x,y) \, dx \, dy$,
        \item Si $f \geq 0$ et que $\iint_D f(x,y) \, dx \, dy = 0$, alors $f$ est nulle sur $D$,
        \item \textbf{Inégalité triangulaire :} $|\iint_D f(x,y) \, dx \, dy| \leq \iint_D |f(x,y)| \, dx \, dy$,
        \item \textbf{Inégalité de Cauchy-Schwarz :} $\left| \iint_D f(x,y) g(x,y) \, dx \, dy \right| \leq \sqrt{\iint_D f(x,y)^2 \, dx \, dy} \sqrt{\iint_D g(x,y)^2 \, dx \, dy}$,
    \end{enumerate}
}

\ndlr{cf. le diable pour la preuve de ce théorème}

\example{
    Intégrons $f(x,y) = \frac{1+x^2}{1+y^2}$ sur le carré $[0,1] \times [0,1]$.\\
    En intégrant d'abord par rapport à $y$, on trouve : $\int_0^1 \frac{1+x^2}{1+y^2} dy = (1+x^2) \int_0^1 \frac{1}{1+y^2} dy = (1+x^2) \left[ \arctan(y) \right]_0^1 = (1+x^2) \frac{\pi}{4}$.\\
    En intégrant ensuite par rapport à $x$, on trouve : $\int_0^1 (1+x^2) \frac{\pi}{4} dx = \frac{\pi}{4} \int_0^1 (1+x^2) dx = \frac{\pi}{4} \left[ x + \frac{x^3}{3} \right]_0^1 = \frac{\pi}{4} \cdot \frac{4}{3} = \frac{\pi}{3}$.\\
    Donc $\iint_{[0,1] \times [0,1]} \frac{1+x^2}{1+y^2} dx dy = \frac{\pi}{3}$.
}

\theorem{Théorème de Fubini}{}{false}{
    Soient $[a,b]$ un segment et $c,d : [a,b] \to \mathbb{R}$ des fonctions continues.\\
    On pose $D = \{ (x,y) \in \mathbb{R}^2 \; | \; x \in [a,b], \; c(x) \leq y \leq d(x) \}$.\\
    Soit $f : D \to \mathbb{R}$ une fonction continue. Alors :
    \[
    \iint_D f(x,y) \, dx \, dy = \int_a^b \left( \int_{c(x)}^{d(x)} f(x,y) \, dy \right) dx
    \]
}

\section{Changement de variables}

\subsection{Matrice jacobienne et jacobien}
\definition{
    Soit $g : \mathbb{R}^2 \to \mathbb{R}^2$ de classe $\mathcal{C}^1$.\\
    La \textbf{matrice jacobienne} de $g$ est la matrice $Jac_g(x,y) = \begin{pmatrix}
        \frac{\partial g_1}{\partial x}(x,y) & \frac{\partial g_1}{\partial y}(x,y) \\
        \frac{\partial g_2}{\partial x}(x,y) & \frac{\partial g_2}{\partial y}(x,y)
    \end{pmatrix}$.
}

\vocabulary{On appelle le \textbf{jacobien} de $g$ le déterminant de la matrice jacobienne de $g$, c'est à dire $det(Jac_g(x,y)) = \frac{\partial g_1}{\partial x}(x,y) \cdot \frac{\partial g_2}{\partial y}(x,y) - \frac{\partial g_1}{\partial y}(x,y) \cdot \frac{\partial g_2}{\partial x}(x,y)$}

\definition{
    Soit $D$ et $\Delta$ deux domaines quarrables de $\mathbb{R}^2$.\\
    Soit $g : D \to \Delta$.\\
    On dit que $g$ est un $\mathcal{C}^1\text{-diffeomorphisme}$ si :
    \begin{itemize}
        \item $g$ est de classe $\mathcal{C}^1$,
        \item $g$ est un bijection de $D$ sur $\Delta$,
        \item $\forall (x,y) \in D, \; det(Jac_g(x,y)) \neq 0$.
    \end{itemize}
}

\theorem{Théorème de changement de variables}{}{false}{
    Soit $D$ et $\Delta$ deux domaines quarrables de $\mathbb{R}^2$.\\
    Soit $g : D \to \Delta$ un $\mathcal{C}^1\text{-diffeomorphisme}$.\\
    Soit $f : \Delta \to \mathbb{R}$ une fonction continue. Alors :
    \[
    \iint_\Delta f(u,v) du dv = \iint_D f(g(x,y)) |det(Jac_g(x,y))| dx dy
    \]
}

\example{
    Calculons l'intégrale de $f(x,y) = x^2 + y^2$ sur le disque unité $\Delta = \{ (x,y) \in \mathbb{R}^2 \mid x^2 + y^2 \leq 1\}$.
    
    On utilise le passage en coordonnées polaires défini par l'application :
    \[ g : D \to \Delta, \quad (r,\theta) \mapsto (r \cos\theta, r \sin\theta) \]
    où le domaine de départ est le rectangle $D = [0,1] \times [0,2\pi]$.
    
    \begin{enumerate}
        \item \textbf{Jacobien :} Le déterminant de la matrice jacobienne est $\det(Jac_g(r,\theta)) = r$. Comme $r \geq 0$ sur $D$, la valeur absolue est inutile ($|r| = r$).
        \item \textbf{Transformation de la fonction :} $f(g(r,\theta)) = (r\cos\theta)^2 + (r\sin\theta)^2 = r^2$.
        \item \textbf{Calcul de l'intégrale :}
    \end{enumerate}
    
    \[
    \iint_\Delta (x^2 + y^2) \, dx \, dy = \iint_D r^2 \cdot r \, dr \, d\theta = \int_0^{2\pi} \left( \int_0^1 r^3 \, dr \right) d\theta
    \]
    
    Par intégrations successives :
    \begin{itemize}
        \item Intégration selon $r$ : $\displaystyle \int_0^1 r^3 \, dr = \left[ \frac{r^4}{4} \right]_0^1 = \frac{1}{4}$.
        \item Intégration selon $\theta$ : $\displaystyle \int_0^{2\pi} \frac{1}{4} \, d\theta = \left[ \frac{\theta}{4} \right]_0^{2\pi} = \frac{2\pi}{4} = \frac{\pi}{2}$.
    \end{itemize}
    
    Finalement : $\displaystyle \iint_\Delta (x^2 + y^2) \, dx \, dy = \frac{\pi}{2}$.
}
\newpage
\tableofcontents

\end{document}